% unit13.tex — R2 unit 13 build driver, The Sourdough Framework: Arabic Edition
%
% Original: The Sourdough Framework by Hendrik Kleinwächter,
% the-sourdough-framework.com, CC BY-SA 4.0. The Arabic edition inherits
% CC BY-SA 4.0.
%
% Unit 13 per the reconstruction plan: ch. 13 storing bread. One table
% (mirrored per the unit 9 RTL table convention, {\textdir TRT} cells),
% two footnotes, two cites, no figures, no flowcharts. Content layer
% re-typeset VERBATIM from the recovery reference
% (recovery-r0/text/13-storing-bread.txt, R0 recovery mining of the
% Phase 23 website build) — no retranslation; the accepted edition
% book-ar-fixed.pdf (b78c2a28…) is the fidelity authority, ch. 13 =
% accepted pp. 148–150. Built with the standing arabic-digits.lua
% (standing md5 a2e3d1daeb3e24ced21873376170cab8) and run.sh per
% Ruling 2 (2026-11-15). Driver scheme per the accepted edition, same
% as units 3–12: NO \frontmatter — chapters are numbered from 1 and
% page numbering runs continuously in Western digits from the cover.
% Unit builds carry natural chapter numbering (ch. 13 builds as ch. 1
% here, the table as 1.1); the R3 full-book assembly regenerates true
% numbering.

\documentclass[paper=a4, twoside=false, fontsize=12pt, parskip=half,
                bibliography=totoc, listof=totoc]{scrbook}
\usepackage{sourdough-ar}
% range-phrase wrapped in \textenglish per the unit 9 dash-font finding
% (2026-11-15), inherited by units 10–13: literal-text ranges
% (\LR{\textenglish{N\textendash M}}) are the accepted convention;
% \numrange with this \textenglish range-phrase is the fallback
% convention. Driver-level change only.
\sisetup{range-phrase={\textenglish{\textendash}}}
\author{Hendrik Kleinw\"achter}
\title{منظومة عمل الخبز الحامض}

\begin{document}
% cover-ar.tex — Arabic cover and title page for The Sourdough Framework
%
% Original: The Sourdough Framework by Hendrik Kleinwächter,
% the-sourdough-framework.com, CC BY-SA 4.0. The Arabic edition inherits
% CC BY-SA 4.0.
%
% Fork of the upstream cover/cover.tex. Page 1 carries the cover raster
% (upstream cover/cover-page.jpg, upstream commit 69636e6) with the Arabic
% overlay of record on the cover band; page 2 is the blank verso; page 3 is
% the typeset title page with the Arabic title, edition date and the
% author/translator credit.
%
% Overlay of record (Phase 26/27, Item 1 applied at R1: 4-word title):
%   Line 1: title, \put(88.45, 7.87), 14/16 Amiri bold, centered
%   Line 2: تأليف: هندريك كلاينفِختر · ترجمة: نائل العيله
%           \put(88.45, 2.92), 8/9 Amiri regular, centered
% The \makeatletter/\if@twoside wrapper is upstream-verbatim (a .tex file;
% the no-\makeatletter rule binds the .sty only).

\thispagestyle{empty}
\setlength{\unitlength}{1mm}
\makeatletter
\if@twoside%
    \noindent\begin{picture}(0,0)(1,-1)
    \put(-16.3,-265){\includegraphics[width=1.33\linewidth]{cover/cover-page.jpg}}
    \end{picture}
\else%
    \noindent\begin{picture}(0,0)(1,-1)
    \put(-23.3,-265){\includegraphics[width=1.33\linewidth]{cover/cover-page.jpg}}
    \end{picture}
\fi%
% Arabic overlay on the cover band (Phase 26/27 of record; 4-word title).
\noindent\begin{picture}(0,0)(1,-1)
    \put(88.45,7.87){\makebox(0,0)[c]{\fontsize{14}{16}\selectfont\bfseries منظومة عمل الخبز الحامض}}%
    \put(88.45,2.92){\makebox(0,0)[c]{\fontsize{8}{9}\selectfont تأليف: هندريك كلاينفِختر · ترجمة: نائل العيله}}%
\end{picture}
\makeatother

\newpage
\thispagestyle{empty}

\rule{1pt}{\textheight} % Vertical line
 % Whitespace between the vertical line and title page text
\hspace{0.05\textwidth}
 % Paragraph box for holding the title page text, adjust the width to move the
% title page left or right on the page
\parbox[b]{0.75\textwidth}{%
{\Huge\bfseries منظومة عمل الخبز الحامض}\\[1\baselineskip] % Title (4-word form, Item 1 of record)
{\large\textenglish{The Sourdough Framework}}\\[2\baselineskip] % Original-title subtitle, accepted p. 2
{\large\textit{الإصدار: 7 نوفمبر 2026}}\\[4\baselineskip] % date frozen to the accepted edition's build date
{\Large هندريك كلاينفِختر \textenglish{(Hendrik Kleinwächter)}}\\[1\baselineskip]
{\large ترجمة: نائل العيله}

% Whitespace between the title block and the copyright text. In the PDF and
% the ebook it holds a note on asking this book: Loafy answers from it. The
% website leaves it out, it has Loafy on every page already. Each place has
% its own via tag, so Loafy can count where readers came from.
\ifdefined\HCode
\vspace{0.5\textheight}
\ifdefined\sdfebook
\textbf{أسئلة أثناء الخَبز؟} اسأل هذا الكتاب. \textenglish{Loafy}، رفيقه المجاني،
يجيب من هذه الصفحات ويسمّي القسم:
\href{https://loafy.ai/loafy?via=tsf-epub-title}{https://loafy.ai}
\fi
\else
\vspace{0.17\textheight}

\noindent
\begin{minipage}[c]{2.4cm}
% Black, not link blue: a light QR code is hard to scan once printed.
% The QR matrix macros must typeset LTR; RTL context corrupts them.
\LRE{\qrcode[nolink,height=2.2cm]{https://loafy.ai/loafy?via=tsf-pdf-qr}}
\end{minipage}%
\hspace{0.4cm}%
\begin{minipage}[c]{\linewidth-2.8cm}
\textbf{أسئلة أثناء الخَبز؟} اسأل هذا الكتاب «لُوفي» \textenglish{(Loafy)}، رفيقه المجاني،
يجيب عنك من صفحاته ويسمّي لك القسم:
\href{https://loafy.ai/loafy?via=tsf-pdf-title}{\nolinkurl{https://loafy.ai}}
\end{minipage}

\vspace{0.17\textheight}
\fi


{\noindent
\begin{flushleft}
    \includegraphics[width=3cm]{cover/CC-BY-SA}\par
الشيفرة المصدرية الكاملة لهذا الكتاب متاحة على
\url{https://github.com/hendricius/the-sourdough-framework/} برخصة
\textenglish{CC-BY-SA}.
راجع \url{https://creativecommons.org/licenses/by-sa/4.0/} لمزيد من التفاصيل.
لا تتردد في الإبلاغ عن الأخطاء أو اقتراح التحسينات. يتوفر إصدار غلاف مقوّى من
الكتاب أيضًا. مزيد من المعلومات هنا:
\url{https://www.breadco.de/hardcover-book}

الإصدار العربي: مبنيّ على الكتاب الأصلي \textenglish{The Sourdough Framework}
لـ هندريك كلاينفِختر، المنشور مرخّصًا برخصة \textenglish{CC BY-SA 4.0}،
ويخضع هذا الإصدار العربي للرخصة نفسها.
\end{flushleft}
}

{\centering\textenglish{Original: The Sourdough Framework by Hendrik
Kleinwächter, the-sourdough-framework.com, CC BY-SA 4.0}\par}
}


{\hypersetup{hidelinks}\tableofcontents}

% storing-bread-ar.tex — Arabic edition of ch. 13, storing bread
%
% Original: The Sourdough Framework by Hendrik Kleinwächter,
% the-sourdough-framework.com, CC BY-SA 4.0. The Arabic edition inherits
% CC BY-SA 4.0.
%
% Phase 28R R2 unit 13 reconstruction. VERBATIM transcription of the R0
% recovery reference (recovery-r0/text/13-storing-bread.txt, mined from
% the Phase 23 website build of the accepted Arabic sources); no text
% retranslated from the English. The accepted edition book-ar-fixed.pdf
% (b78c2a28b574288801d8000a7249af3a) is the fidelity authority; ch. 13 =
% accepted pp. 148–150. Short caption (خيارات تخزين الخبز) recovered
% from the accepted edition's list-of-tables page (قائمة الجداول,
% entry 13.1, accepted p. 192).
%
% No figures, no flowcharts, no recipes in this chapter. One table
% (upstream [!htb] spec kept verbatim; placement verified against the
% accepted render — table on the chapter opening page under the intro,
% caption at the page bottom). Two footnotes, two cites (upstream keys
% storing+bread, freezing+toasting+bread).
%
% Direction treatments per the standing conventions:
%  - the table is mirrored per the unit 9 RTL table convention
%    (columns visual LTR, label column الطريقة rightmost); every
%    non-empty cell wrapped {\textdir TRT ...} (unit 9, BINDING);
%  - Latin terms ride \textenglish{} islands: (breadcrumbs), Knödel,
%    (mold), (airtight container), (ziplock bag), (freezer), (thawing),
%    (toaster) — paren-font residual (Pagella vs Amiri parens)
%    documented, R3 harmonization agenda;
%  - digit ranges as literal-text LTR islands per the unit 9 dash-font
%    convention: \LR{\textenglish{1\textendash 2}} and
%    \LR{\textenglish{3\textendash 4}};
%  - bare digit runs ride per-token \LR islands: \LR{10} إلى \LR{15}
%    ثانية, عام \LR{2008};
%  - the single percentage rides \LR{\qty{40}{\percent}} (unit 11
%    convention);
%  - cites ride ~\cite{...} with the upstream keys (units 9–12
%    pattern);
%  - footnote markers AFTER the sentence period (accepted p. 149
%    token order 'طرية.1' and '.2' verified in the render), unlike the
%    upstream English (marker before period);
%  - the same-unit table ref renders \LR{\ref{...}} per the same-unit
%    ref convention (unit build shows 1.1 where the accepted shows
%    13.1 — unit-numbering artifact, R3 assembly regenerates true
%    numbering);
%  - the accepted edition's own text-layer tatweel in 'فالـكيس' /
%    'ﺍلـكيس' (p. 149) is a justification kashida artifact of the
%    accepted toolchain, not source text — the recovery reference (mined
%    from the accepted sources) has فالكيس / الكيس without tatweel, and
%    the transcription follows the recovery reference.

\chapter{تخزين الخبز}\label{ch:storing-bread}

\begin{quoting}
يناقش هذا الفصل طرقًا مختلفة لتخزين خبزك، ولكل منها مزاياها وعيوبها. وبهذا يمكن الاستمتاع بخبزك في وقت لاحق بأفضل صورة ممكنة.
\end{quoting}

تجد ملخصًا في الجدول~\LR{\ref{table:bread-storage}}، مع التفاصيل والشرح في بقية هذا الفصل.
\begin{table}[!htb]
    \centering
        % table-storing-bread-overview-ar.tex — Arabic edition of the bread
% storage options table (ch. 13, جدول 13.1)
%
% Original: The Sourdough Framework by Hendrik Kleinwächter,
% the-sourdough-framework.com, CC BY-SA 4.0. The Arabic edition inherits
% CC BY-SA 4.0.
%
% Phase 28R R2 unit 13 reconstruction. Column geometry verbatim from the
% upstream source tables/table-storing-bread-overview.tex
% ({@{}p{0.23\textwidth}p{0.33\textwidth}p{0.33\textwidth}@{}} with
% booktabs rules and \arraystretch 1.6), mirrored per the standing RTL
% table convention: columns visual LTR, the label column (الطريقة)
% rightmost, the two wide description columns (العيوب leftmost,
% المزايا middle). Cell text recovered verbatim from the accepted
% edition's render of p. 148 (book-ar-fixed.pdf,
% b78c2a28b574288801d8000a7249af3a) via the R0 recovery reference
% (recovery-r0/text/13-storing-bread.txt) — the text layer extracts
% cleanly, so the cells are recovered against the accepted text layer,
% not retranslated.
%
% Every non-empty cell is wrapped {\textdir TRT ...} per the unit 9
% cell-direction convention (BINDING). No Latin or digit runs in this
% table, so no \LR islands are needed.

\def\arraystretch{1.6}
\begin{tabular}{@{}p{0.33\textwidth}p{0.33\textwidth}p{0.23\textwidth}@{}}
\toprule
\thead{{\textdir TRT العيوب}} & \thead{{\textdir TRT المزايا}} & \thead{{\textdir TRT الطريقة}} \\
\midrule
{\textdir TRT يجفّ الخبز سريعًا جدًّا.} & {\textdir TRT الأسهل على الإطلاق. الأفضل للخبز الذي يُؤكل خلال يوم واحد. تبقى القشرة عادةً مقرمشة ما لم تكن الرطوبة مرتفعة جدًّا.} & {\textdir TRT درجة حرارة الغرفة} \\
{\textdir TRT يتعين تحميص الخبز لتعود القشرة إلى مقرمشتها. يتعرّض للعَفَن بسرعة أكبر.} & {\textdir TRT جيد حتى أسبوع.} & {\textdir TRT درجة حرارة الغرفة داخل وعاء محكم الإغلاق} \\
{\textdir TRT يتعين تحميص الخبز. يتطلب ثلاجة وطاقة.} & {\textdir TRT يبقى الخبز جيدًا لأسابيع. وقد يجفّ قليلًا عند عدم استخدام وعاء محكم الإغلاق.} & {\textdir TRT الثلاجة} \\
{\textdir TRT يتطلب فكّ التجمّد ثم التحميص. يتطلب مُجمِّدًا وطاقة.} & {\textdir TRT يبقى الخبز جيدًا لسنوات.} & {\textdir TRT المُجمِّد} \\
\bottomrule
\end{tabular}

        \caption[خيارات تخزين الخبز]{جدول يوضح مزايا وعيوب خيارات تخزين الخبز المختلفة.}%
        \label{table:bread-storage}
\end{table}

\section{درجة حرارة الغرفة}

الطريقة الأكثر شيوعًا هي تخزين خبزك في درجة حرارة الغرفة. وبعد أن تأخذ شريحة من الخبز، خزّن خبزك مع لُبّه باتجاه الأسفل.

تعمل هذه الطريقة بشكل رائع إذا أردت أكل خبزك خلال يوم واحد. فتبقى القشرة مقرمشة ولا تصبح طرية.\footnote{كلما ارتفعت الرطوبة في غرفتك أصبحت القشرة طرية بسرعة أكبر.} أما أكبر عيوب هذه الطريقة فهو أن الخبز يتصلّب بسرعة. فمع مضي الوقت يتبخر من لُبّ عجينك ماء أكثر فأكثر. وفي النهاية سيصبح الخبز صلبًا جدًّا يستحيل أكله. وكلما زاد الماء الذي تستخدمه في صنع الخبز بقي الخبز جيدًا مدة أطول. فالوصفة منخفضة الترطيب قد تجفّ بعد \LR{\textenglish{1\textendash 2}} يوم، بينما يحتاج الخبز مرتفع الترطيب إلى \LR{\textenglish{3\textendash 4}} أيام ليجفّ.

وبمجرد أن يجفّ خبزك يمكنك تمريره تحت ماء الصنبور مدة \LR{10} إلى \LR{15} ثانية تقريبًا. يتيح هذا الحمّام المائي للنشويات الموجودة في اللُبّ أن تمتص كمية كبيرة من الماء. ثم اخبز خبزك مجددًا في الفرن. سيكون الرغيف الناتج جيدًا شبه كأنه جديد.

وهناك خيار آخر للخبز اليابس وهو استخدامه في صنع البقسماط \textenglish{(breadcrumbs)}. ويمكن خلط هذه الفتاتات الخبزية في أرغفة لاحقة. كما يمكن استخدامها مكونات أساسية في وصفات أخرى مثل \textenglish{Knödel}.\footnote{\textenglish{Knödel} هو طبق نمساوي يقوم على الخبز القديم. تُخلط الفتاتات وخبز اليوم السابق مع البيض، وتُضاف أحيانًا السبانخ أو اللحم المقدد. ثم يُسلق الخليط في ماء مملح.}

\section{درجة حرارة الغرفة داخل وعاء}

كما في الخيار السابق، يمكنك أيضًا تخزين خبزك داخل وعاء. وقد يكون هذا كيسًا ورقيًّا أو كيسًا بلاستيكيًّا أو صندوق تخزين خبز. فالكيس الورقي ومعظم صناديق الخبز ليست محكمة الإغلاق تمامًا، مما يسمح ببعض الهواء بالانتشار خارج الوعاء. وهذا يعني أيضًا أن الخبز سيجفّ قليلًا.

وعند استخدام كيس محكم مثل الكيس البلاستيكي سيحتفظ الخبز بكمية كبيرة من الرطوبة. فيبقى الخبز جيدًا مدة أطول. غير أنه في الوقت نفسه ستفقد القشرة قرمشتها أيضًا. فبعض الماء ينتشر في الكيس ثم يعيد القشرةُ امتصاصَه. وإذا أردت القشرة المقرمشة فالخيار الأفضل هو تحميص خبزك.

ومن المشكلات الأخرى لأوعية التخزين تلوث العَفَن \textenglish{(mold)} الطبيعي. ففي اللحظة التي يخرج فيها خبزك من الفرن تبدأ بتلقّح بأبواغ العَفَن المحمولة في الهواء. وهذه الأبواغ صغيرة جدًّا لدرجة المجهر وهي في كل مكان. وتنمو أبواغ العَفَن بأفضل صورة في بيئة رطبة. وبوضع عجينك داخل وعاء تكون قد صنعت جنةً للعَفَن. فالعجين العادي القائم على الخميرة سيبدأ بالعَفَن خلال أيام قليلة هكذا. أما الخبز القائم على البادئة الحامضة فيبقى جيدًا مدة أطول لأن الحموضة مانع طبيعي للعَفَن.

\section{الثلاجة}

في تجربتي الخاصة يعمل تخزين الخبز داخل الثلاجة بشكل جيد طالما تستخدم وعاءً محكم الإغلاق \textenglish{(airtight container)}، حتى إن بعض المصادر تقول إن الخبز يجفّ داخل الثلاجة~\cite{storing+bread}. فمن المفترض أن الثلاجة تشجع السائل الموجود في اللُبّ على الانتقال إلى سطح الخبز.

لكن في تجربتي الحيلة هي استخدام وعاء قابل للإغلاق. فمع كيس بسحّاب محكم \textenglish{(ziplock bag)} ستبقى الرطوبة الزائدة في الكيس وتضمن ألا يجفّ الخبز بسرعة. وفي درجة حرارة الغرفة كان هذا سيسبب عَفَن خبزك. وعند درجات الحرارة المنخفضة يمكن أن يبقى الخبز جيدًا هكذا لأسابيع. غير أن القشرة ستفقد قرمشتها، ولهذا يُنصح بالتحميص.

\section{التجميد}

وهناك خيار رائع آخر للتخزين طويل الأمد وهو استخدام مُجمِّدك \textenglish{(freezer)}. قطّع الرغيف كله وأنشئ أجزاءً يمكنك استهلاكها خلال يوم واحد. وخزّن كل جزء في وعاء منفصل وضعها داخل مُجمِّدك.

وعندما تريد أكل خبز طازج، افتح أحد الأوعية في الصباح واترك الخبز يفكّ تجمّده \textenglish{(thawing)} على مدى ساعات قليلة. هذا ضروري كي تتمكن من فصل الشرائح المتجمدة المتلاصقة بسهولة. ثم حمّص الشرائح في محمصة الخبز \textenglish{(toaster)} أو اخبزها في الفرن حتى تبلغ القرمشة التي تعجبك.

وهذا الخيار رائع للتخزين طويل الأمد جدًّا. وأنا شخصيًّا أحب أن تكون لديّ بضع شرائح خبز مجمدة بوصفها احتياطيًّا طارئًا حين لا يكون لديّ وقت للخَبز.

وتلمح دراسة أجريت عام \LR{2008} إلى احتمال وجود فوائد صحية لتجميد خبزك وتحميصه. فبفعل ذلك قد تصبح جزيئات النشا أكثر مقاومة للهضم، مما يخفض استجابة سكر الدم في جسمك بما يقارب \LR{\qty{40}{\percent}}~\cite{freezing+toasting+bread}.


\end{document}
